
Par
ailleurs, on a @@M0@@, où @@M1@@ parcourt l'ensemble des éléments non nuls de @@M2@@. Par suite, @@M3@@, et :
@@M4@@
d'après EGA IV 8.2.3 et 8.2.5. Comme les schémas envisagés sont noethériens, et les morphismes de transition affines et plats, on a d'après IV 3.2.3 :
@@M5@@
Comme les homomorphismes «~image inverse~» sur les @@M6@@ sont linéaires par rapport aux @@M7@@ (IV (2.12.2)), il existe un ouvert non vide @@M8@@ de @@M9@@ tel que dans @@M10@@ on ait :
@@M11@@
Soit @@M12@@ un sous-schéma fermé de @@M13@@ de support @@M14@@. Alors @@M15@@ est surjective d'après l'hypothèse de récurrence, et par suite le diagramme (1.2.1) montre que :
@@M16@@
donc, comme @@M17@@, @@M18@@ est surjective.

\textbf{Corollaire 1.3.} \uline{Soit} @@M19@@ \uline{un morphisme plat de schémas noethériens}. \uline{On suppose que pour tout point} @@M20@@ \uline{de} @@M21@@, \uline{l'homomorphisme} :
@@M22@@
\uline{est surjectif}. \uline{Alors l'homomorphisme déduit de} @@M23@@ :
@@M24@@
\uline{est surjectif}.

